% AUTOMATISCH ERZEUGT aus dem SysML-Modell mit tools/generate_docs.py.
% Nicht von Hand aendern: Inhalte im Modell pflegen, Layout im Kapitel.
DEC-001 & Das Zielsystem ist ein Indoor-Demonstrator. & Angenommen \\
DEC-002 & Einsatzort ist ein produktionsaehnlicher Schulbereich. & Angenommen \\
DEC-003 & Transportgut ist ein Werkstuecktraeger eines Transfersystems. & Angenommen \\
DEC-004 & Die Beladung erfolgt zunaechst extern durch Mensch oder Transfersystem. & Angenommen \\
DEC-005 & Der Fahrbereich umfasst mehrere Raeume und einen Flur. & Angenommen \\
DEC-006 & Offene Tueren muessen befahrbar sein. & Angenommen \\
DEC-007 & Die MVP-Geschwindigkeit liegt bei ca. $0{,}5\,\mathrm{m/s}$. & Angenommen \\
DEC-008 & Die Zielbetriebsdauer der ersten Version liegt bei ca. 2 h. & Angenommen \\
DEC-009 & Transportauftraege werden zunaechst ueber eine Web-App ausgeloest. & Angenommen \\
DEC-010 & Eine einfache Bedienoberflaeche fuer Status, Stopp und Rueckkehr zur Ladestation ist vorgesehen. & Angenommen \\
DEC-011 & Rueckkehr zur Ladestation gehoert zum Zielsystem. & Angenommen \\
DEC-012 & Hoehere Geschwindigkeit wird als spaetere Erweiterung gefuehrt. & Angenommen \\
DEC-013 & Ein betrieblicher Stopp pausiert den laufenden Transportauftrag. Auftrags-ID, Start und Ziel bleiben erhalten. Fortsetzung nur nach ausdruecklicher manueller Freigabe; Sicherheits- und Routenpruefung werden dabei nicht umgangen. & Angenommen (28.10.2026) \\
DEC-014 & Verifikation erfolgt zunaechst ausschliesslich in der Simulation, ohne reale Hardware. Der Simulator ist noch nicht ausgewaehlt. & Angenommen (28.10.2026) \\
DEC-015 & Ein pausierter Auftrag kann vom Bediener abgebrochen werden; er wird verworfen, der Roboter meldet gestoppt. Aus gestoppt, abgeschlossen, abgelehnt, blockiert, Ladebereich erreicht, Fehler und sicherem Zustand fuehrt nur eine manuelle Quittierung zurueck in den Bereitschaftszustand. Aus Fehler nur bei behobenem Fehler, aus sicherem Zustand nur wenn die Bewegung wieder sicher ist. Die Quittierung setzt keinen Not-Halt zurueck. & Angenommen (05.11.2026) \\
DEC-016 & Unsichere Situationen fuehren in den sicheren Zustand, technische Fehler in den Fehlerzustand, und zwar in jedem Betriebszustand ausser ausgeschaltet, Fehlerzustand und sicherem Zustand. Das gilt auch im Stillstand, beim Hochfahren und an der Ladestation; neue Zustaende sind automatisch eingeschlossen. & Angenommen (05.11.2026) \\
DEC-017 & Erreicht der Ladezustand waehrend eines Auftrags den Schwellwert, meldet der Roboter eine Warnung im Missionsstatus und faehrt weiter. Er unterbricht den Auftrag nicht selbst und faehrt nicht selbst zur Ladestation; der Bediener entscheidet. Ein pausierter Auftrag kann zum Laden geschickt werden und bleibt dabei erhalten; nach dem Laden setzt der Bediener ihn manuell fort oder verwirft ihn per Quittierung. & Angenommen (05.11.2026) \\
DEC-018 & Ein betrieblicher Stopp waehrend der Rueckkehr zur Ladestation haelt den Roboter an und fuehrt in den eigenen Zustand Rueckkehr pausiert. Fortsetzen mit manueller Freigabe setzt die Rueckkehr fort, Quittierung bricht sie ab. Ein erhaltener Auftrag bleibt bis zur Quittierung erhalten. & Angenommen (05.11.2026) \\
